\subsection{GridWorld}

GridWorld est un environnement de navigation sur grille 2D où l'agent doit atteindre un objectif tout en évitant les obstacles et les zones dangereuses.

\begin{itemize}
    \item[] \textbf{État :} vecteur binaire de dimension 75 (grille 5×5 avec 3 canaux) :
    \begin{itemize}
        \item Canal 0 : position de l'agent,
        \item Canal 1 : position de l'objectif,
        \item Canal 2 : positions des obstacles et zones dangereuses.
    \end{itemize}

    \item[] \textbf{Actions :} quatre actions discrètes :
    \begin{enumerate}
        \item déplacer vers le haut,
        \item déplacer vers la droite,
        \item déplacer vers le bas,
        \item déplacer vers la gauche.
    \end{enumerate}

    \item[] \textbf{Récompense :}
    \begin{itemize}
        \item $+1$ en atteignant l'objectif,
        \item $-1$ en entrant dans une zone dangereuse,
        \item $-0.001$ à chaque pas de temps (coût de déplacement).
    \end{itemize}

    \item[] \textbf{Terminalité :} l'épisode s'arrête si :
    \begin{itemize}
        \item[] l'agent atteint l'objectif,
        \item[] l'agent entre dans une zone dangereuse,
        \item[] un nombre maximum d'étapes est atteint.
    \end{itemize}
\end{itemize}

%-----------------------------------------
\subsubsection{Architecture du Réseau de Neurones}

L'extraction de features utilise \texttt{MapFromDumpExtractor2} qui encode les positions de l'agent, de l'objectif et des obstacles dans une représentation binaire de dimension 75.

Le réseau $\hat{Q}(s,a;\theta)$ est un perceptron multicouche (MLP) avec l'architecture suivante :

\[
\text{État} \xrightarrow{75 \text{ inputs}} \text{Couche } 1 \xrightarrow{30 \text{ neurones}} \text{Couche } 2 \xrightarrow{30 \text{ neurones}} \text{Q-values} \xrightarrow{4 \text{ outputs}}
\]

\textbf{Architecture implémentée :}

\begin{table}[h!]
    \centering
    \begin{tabular}{|l|c|c|c|}
    \hline
    \textbf{Couche} & \textbf{Entrée} & \textbf{Sortie} & \textbf{Activation} \\
    \hline
    Input (État) & - & 75 & - \\
    \hline
    Dense (Cachée 1) & 75 & 30 & $\tanh$ \\
    \hline
    Dense (Cachée 2) & 30 & 30 & $\tanh$ \\
    \hline
    Output (Q-values) & 30 & 4 & Linéaire \\
    \hline
    \end{tabular}
    \caption{Architecture du réseau Q pour GridWorld-v0 : 75 $\to$ 30 $\to$ 30 $\to$ 4}
\end{table}

\textbf{Justification :} GridWorld possède un espace d'état beaucoup plus complexe (75D). L'architecture à \textbf{2 couches cachées de 30 neurones} chacune permet de capturer les dépendances spatiales et de naviguer efficacement cette grille. Les deux couches $\tanh$ introduisent la non-linéarité suffisante pour apprendre une politique de navigation robuste.

%-----------------------------------------
\subsubsection{Hyperparamètres de l'Entraînement}

\begin{table}[h!]
    \centering
    \small
    \begin{tabular}{|l|c|l|}
    \hline
    \textbf{Paramètre} & \textbf{Valeur} & \textbf{Justification} \\
    \hline
    Learning rate ($\alpha$) & $0.0001$ & Apprentissage lent pour grande capacité réseau \\
    \hline
    Discount factor ($\gamma$) & $0.99$ & Horizon court relativement, proximité de l'objectif \\
    \hline
    Initial $\varepsilon$ & $0.1$ & 10\% d'exploration aléatoire \\
    \hline
    $\varepsilon$ decay & $0.9999$ & Décroissance modérée sur 100k épisodes \\
    \hline
    Minimum $\varepsilon$ & $0.01$ & 1\% minimum pour exploration \\
    \hline
    Batch size & $100$ & Taille standard \\
    \hline
    Buffer capacity & $10,000$ & Pour problème de moyenne complexité \\
    \hline
    Update frequency & $10$ & Chaque 10 événements \\
    \hline
    Target network update & $1,000$ & Stabilité pour apprentissage complexe \\
    \hline
    \end{tabular}
\end{table}

%-----------------------------------------
\subsubsection{Résultats obtenus}

\begin{figure}[h!]
    \centering
    \begin{subfigure}{0.45\linewidth}
        \centering
        \includegraphics[width=\linewidth]{PLACEHOLDER_gridworld_reward_train.png}
        \caption{Récompense dans l'entraînement}
    \end{subfigure}
    \hfill
    \begin{subfigure}{0.32\linewidth}
        \centering
        \includegraphics[width=\linewidth]{PLACEHOLDER_gridworld_qvalue.png}
        \caption{Q-valeur moyenne}
    \end{subfigure}
    \hfill
    \begin{subfigure}{0.32\linewidth}
        \centering
        \includegraphics[width=\linewidth]{PLACEHOLDER_gridworld_reward_test.png}
        \caption{Récompense dans le test}
    \end{subfigure}

    \caption{Performance des trois métriques pour les variantes GridWorld}
\end{figure}

\subsubsection*{Interprétation des résultats}

[À compléter après exécution de compare\_variants\_gridworld.py]

%-----------------------------------------
%-----------------------------------------
%-----------------------------------------

\subsection{LunarLander}

LunarLander est un problème classique où l'agent doit atterrir un vaisseau spatial de manière sûre sur la lune en controlant la position et la rotation.

\begin{itemize}
    \item[] \textbf{État :} vecteur réel de dimension 8 :
    \begin{enumerate}
        \item position horizontale du vaisseau ($x$),
        \item position verticale du vaisseau ($y$),
        \item vitesse horizontale ($v_x$),
        \item vitesse verticale ($v_y$),
        \item angle de rotation ($\theta$),
        \item vitesse angulaire ($\omega$),
        \item contact jambe gauche (0 ou 1),
        \item contact jambe droite (0 ou 1).
    \end{enumerate}

    \item[] \textbf{Actions :} quatre actions discrètes :
    \begin{enumerate}
        \item ne rien faire,
        \item allumer moteur principal,
        \item allumer moteur latéral gauche,
        \item allumer moteur latéral droit.
    \end{enumerate}

    \item[] \textbf{Récompense :}
    \begin{itemize}
        \item $+100$ pour un atterrissage réussi,
        \item $-100$ en cas de crash,
        \item récompense basée sur la position et la vitesse (-1.5 $\times$ position $y$).
    \end{itemize}

    \item[] \textbf{Terminalité :} l'épisode s'arrête si :
    \begin{itemize}
        \item[] le vaisseau touche le sol (atterrissage réussi ou crash),
        \item[] un nombre maximum d'étapes est atteint.
    \end{itemize}
\end{itemize}

%-----------------------------------------
\subsubsection{Architecture du Réseau de Neurones}

Le réseau $\hat{Q}(s,a;\theta)$ est un perceptron multicouche (MLP) avec l'architecture suivante :

\[
\text{État} \xrightarrow{8 \text{ inputs}} \text{Couche cachée} \xrightarrow{\textbf{200 neurones}} \text{Q-values} \xrightarrow{4 \text{ outputs}}
\]

\textbf{Architecture implémentée :}

\begin{table}[h!]
    \centering
    \begin{tabular}{|l|c|c|c|}
    \hline
    \textbf{Couche} & \textbf{Entrée} & \textbf{Sortie} & \textbf{Activation} \\
    \hline
    Input (État) & - & 8 & - \\
    \hline
    Dense (Cachée) & 8 & 200 & $\tanh$ \\
    \hline
    Output (Q-values) & 200 & 4 & Linéaire \\
    \hline
    \end{tabular}
    \caption{Architecture du réseau Q pour LunarLander-v2 : 8 $\to$ 200 $\to$ 4}
\end{table}

\textbf{Justification :} LunarLander possède un espace d'état continu relativement petit (8D), mais requiert une approximation très précise de la fonction Q pour piloter correctement le vaisseau. La couche cachée de \textbf{200 neurones} avec activation $\tanh$ fournit la capacité d'approximation nécessaire pour capturer les dynamiques complexes du contrôle aérien.

%-----------------------------------------
\subsubsection{Hyperparamètres de l'Entraînement}

\begin{table}[h!]
    \centering
    \small
    \begin{tabular}{|l|c|l|}
    \hline
    \textbf{Paramètre} & \textbf{Valeur} & \textbf{Justification} \\
    \hline
    Learning rate ($\alpha$) & $0.0001$ & Convergence lente pour problème difficile \\
    \hline
    Discount factor ($\gamma$) & $0.99$ & Récompenses futures importantes, horizon moyen \\
    \hline
    Initial $\varepsilon$ & $0.1$ & 10\% d'exploration aléatoire \\
    \hline
    $\varepsilon$ decay & $0.99999$ & Décroissance très lente, exploration prolongée \\
    \hline
    Minimum $\varepsilon$ & $0.01$ & 1\% minimum pour exploration \\
    \hline
    Batch size & $100$ & Taille standard \\
    \hline
    Buffer capacity & $100,000$ & Très grande pour problème complexe et stochastique \\
    \hline
    Update frequency & $10$ & Chaque 10 événements \\
    \hline
    Target network update & $1,000$ & Stabilité critique pour ce problème \\
    \hline
    \end{tabular}
\end{table}

%-----------------------------------------
\subsubsection{Résultats obtenus}

\begin{figure}[h!]
    \centering
    \begin{subfigure}{0.45\linewidth}
        \centering
        \includegraphics[width=\linewidth]{PLACEHOLDER_lunar_reward_train.png}
        \caption{Récompense dans l'entraînement}
    \end{subfigure}
    \hfill
    \begin{subfigure}{0.32\linewidth}
        \centering
        \includegraphics[width=\linewidth]{PLACEHOLDER_lunar_qvalue.png}
        \caption{Q-valeur moyenne}
    \end{subfigure}
    \hfill
    \begin{subfigure}{0.32\linewidth}
        \centering
        \includegraphics[width=\linewidth]{PLACEHOLDER_lunar_reward_test.png}
        \caption{Récompense dans le test}
    \end{subfigure}

    \caption{Performance des trois métriques pour les variantes LunarLander}
\end{figure}

\subsubsection*{Interprétation des résultats}

[À compléter après exécution de compare\_variants\_lunar.py]
